\chapter{Identify the dot and cross products}

The ordinary three-dimensional dot product is
\[
\mathbf P\cdot\mathbf Q
=P_1Q_1+P_2Q_2+P_3Q_3.
\]

Therefore,
\[
\operatorname{Sc}(PQ)
=P_0Q_0-\mathbf P\cdot\mathbf Q.
\]

The ordinary three-dimensional cross product is
\[
\mathbf P\times\mathbf Q
=
\begin{pmatrix}
P_2Q_3-P_3Q_2\\
P_3Q_1-P_1Q_3\\
P_1Q_2-P_2Q_1
\end{pmatrix}.
\]

The vector part can consequently be written as
\[
\operatorname{Vec}(PQ)
=P_0\mathbf Q+Q_0\mathbf P+\mathbf P\times\mathbf Q.
\]

Thus the complete quaternion product is
\[
\boxed{
PQ=
\left(P_0Q_0-\mathbf P\cdot\mathbf Q\right)
+
\left(P_0\mathbf Q+Q_0\mathbf P+\mathbf P\times\mathbf Q\right)
}.
\]

Here the first parentheses contain the scalar part, and the second parentheses contain the vector part.

\section*{4. Why the cross product appears}

For example, the terms contributing to the \(i\)-component include
\[
(jP_2)(kQ_3)=iP_2Q_3,
\]
whereas
\[
(kP_3)(jQ_2)=-iP_3Q_2.
\]

Together they give
\[
i(P_2Q_3-P_3Q_2),
\]
which is the \(i\)-component of
\(\mathbf P\times\mathbf Q\).

The \(j\)- and \(k\)-components arise in the same way:
\[
j(P_3Q_1-P_1Q_3),
\]
\[
k(P_1Q_2-P_2Q_1).
\]

The cross product therefore emerges from the antisymmetric part of quaternion multiplication.
